\documentclass[12pt]{article}

%Lei de Gauss para uma cavidade cilindrica
%Faro, Mon Mar  3 01:59:44 WET 2014
%Written by Tordar

\usepackage{graphicx}
\usepackage{pst-all}

\begin{document}

\pagestyle{empty}

\begin{pspicture}(-5,-5)(5,5)
\psframe[linewidth=1.5mm,linecolor=white,fillstyle=vlines](-5,-4)(5,4)
\pscircle[linewidth=1.5mm,linecolor=black,fillstyle=solid,fillcolor=white](1,0){3}
\rput{0}(0,2){\scalebox{5}{$\cdot$}}
\psline[linewidth=1mm,linecolor=blue,arrows=->](-4,0)(1,0)
\psline[linewidth=1mm,linecolor=blue,arrows=->]( 1,0)(0,2)
\psline[linewidth=1mm,linecolor=blue,arrows=->](-4,0)(0,2)
\rput{0}(-2.5,-0.5){\scalebox{1.5}{\psframebox*{$\vd$}}}
\rput{0}(-2.5, 1.5){\scalebox{1.5}{\psframebox*{$\mathbf{r}$}}}
\rput{0}(1.2,1){\scalebox{1.5}{\psframebox*{$\mathbf{r}'$}}}
\rput{0}(0.5,2){\scalebox{1.2}{$P$}}
\end{pspicture}

\end{document}
